\documentclass[11pt]{article}
\usepackage[margin=1in]{geometry}
\usepackage{iftex}
\ifPDFTeX
  \usepackage[T1]{fontenc}
  \usepackage[utf8]{inputenc}
\fi
\usepackage{amsmath,amssymb}
\usepackage{graphicx}
\usepackage{booktabs}
\usepackage{longtable}
\usepackage{etoolbox}
\AtBeginEnvironment{longtable}{\small}
\AtBeginEnvironment{figure}{\small}
% figure captions: small, bold label, inset from the margins, with clear space before the body text
\makeatletter
\long\def\@makecaption#1#2{%
  \vskip\abovecaptionskip
  {\small\leftskip=1.5em\rightskip=1.5em
   \sbox\@tempboxa{\textbf{#1.}\ #2}%
   \ifdim\wd\@tempboxa>\dimexpr\hsize-3em\relax
     \noindent\textbf{#1.}\ #2\par
   \else
     \global\@minipagefalse\hb@xt@\hsize{\hfil\box\@tempboxa\hfil}%
   \fi}%
  \vskip\belowcaptionskip}
\makeatother
\setlength{\textfloatsep}{18pt plus 4pt minus 2pt}
\setlength{\intextsep}{16pt plus 4pt minus 2pt}
\usepackage{array}
\usepackage{calc}
% caption package removed 2026-09-16: it breaks pandoc 3 uncaptioned longtables (\LTcaptype{none})
\usepackage{microtype}
\usepackage{xcolor}
\usepackage[numbers,sort&compress]{natbib}
\usepackage{hyperref}
\hypersetup{colorlinks=true,linkcolor=blue!60!black,citecolor=blue!60!black,urlcolor=blue!60!black}
\providecommand{\tightlist}{\setlength{\itemsep}{0pt}\setlength{\parskip}{0pt}}
\providecommand{\pandocbounded}[1]{#1}
\setlength{\parskip}{4pt}
\setlength{\parindent}{0pt}
\graphicspath{{figures/}}
\title{Understanding through Perturbation: lowering classifier-free guidance as
a dose-response assay for two text-to-image latent models}
\author{Youssef Hassan\\ \small YWSF Lab (independent
researcher)\\ \small \texttt{youssefhassan13@gmail.com}}
\date{October 2026}
\begin{document}
\maketitle
\begin{abstract}
I probe a generative model by lowering one control, pre-registered, and scoring
what breaks with a rubric borrowed from Klüver. I generated 860 images on two
latent text-to-image checkpoints, SDXL and SD 3.5: 6 prompts, 7 guidance values,
10 seeds, scored blind by two vision-language judges on criteria I committed to
git before the run; two humans rated a subset. Lower classifier-free guidance
was associated with greater rubric-scored visual breakdown, and the association
remained after adjustment for two no-reference quality proxies, robustly to
prompt resampling on SDXL but not on SD 3.5. SDXL met every pre-registered gate
under the post-data amended judge panel. SD 3.5 missed the slope gate: its slope
of -0.182 sits above the -0.20 line, and its interval straddles that line, so an
effect of the registered size is neither shown nor ruled out. Looking
afterwards, almost all of SD 3.5's drop happens between g = 1 and g = 2. A
second pre-registered rubric, adapted from Suzuki's axes, replicated on SDXL and
failed on SD 3.5, where veridicality moved opposite to the registered direction.
Qwen2.5-VL-7B, the original second judge, failed silently, returning well-formed
zeros on visibly broken images, and was caught only by a human subset after the
run; Qwen3-VL-8B failed the same way in screening. Replacing the 7B with
Qwen3-VL-32B, a post-data amendment, raised composite kappa from 0.290 to 0.562
(SDXL) and 0.158 to 0.440 (SD 3.5; 0.394 without the empty-prompt baselines);
the two humans agreed at 0.567, and with the best judge at 0.337 and 0.426. An
earlier experiment found no Klüver form constants, but its judge also flagged
them on 45\% of ordinary images, above a registered ceiling of 20\%, so that
null is uninterpretable under its own rule. The psychedelic analogy supplies the
rubric and is not tested; code, pre-registrations, images and every judge reply
are public.
\end{abstract}
\section{Introduction}\label{introduction}

I set out to understand a generative system by breaking it in a controlled way.
Klüver catalogued what human vision does when it breaks under mescaline,
reducing the first stages of the intoxication to a small set of form constants
\citep{kluver1966}. The four fields I score, each 0 to 3, are reduplication (too
many copies), fragmentation (broken pieces), condensation (fused objects) and
distortion (warped shapes); they are borrowed from that inventory; the second
rubric is adapted from Suzuki's axes \citep{suzuki2024}. The adaptation is not
cosmetic: spontaneity here means content the prompt did not ask for, where
Suzuki's means independence from sensory input, and veridicality here combines
realism and coherence. The knob I turn is classifier-free guidance, g: how hard
the sampler is pushed toward the prompt, where at g = 1 the prompt is used but
not amplified. It is a control most text-to-image pipelines expose
(guidance-distilled models do not), and turning it down is cheap, reversible and
pre-registrable. The original question was whether diffusion output breaks down
into Klüver's shapes, the lattices, tunnels and spirals that neural-field theory
links to the organisation of early visual cortex \citep{bressloff2001}. Exp 01
returned an apparent null. A later registered check (Exp 02) showed that its
judge did find form constants that were really there, but also reported them in
almost half of ordinary images, so its silence on diffusion output cannot be
read as their absence. This paper is about the judge-scored object-level
breakdown that appeared instead, and the instrument needed to see it.

Classifier-free guidance combines a conditional and an unconditional model
prediction at each step, weighted so that more guidance buys fidelity and costs
diversity \citep{ho2022cfg}. SDXL predicts noise and SD 3.5 predicts a velocity,
so ``prediction'' and not ``score'' is the term that covers both. Ho and
Salimans write the tempered distribution \(p(x \mid c)\,p(c \mid x)^{w}\) for
classifier guidance, and present classifier-free guidance as inspired by an
implicit classifier rather than equal to it \citep{ho2022cfg}; more bluntly, it
sharpens sampling toward high-likelihood regions of the conditional
\citep{kynkaanniemi2024interval, karras2024autoguidance}. Lowering g toward 1
removes the amplification of the prompt and approaches the unguided conditional
model, which in the diffusers convention used here is g = 1; only below 1 does
the sampler move toward the unconditional model. Which side of the
psychedelic-neuroscience analogy that corresponds to is a modelling choice; both
readings are defensible. Read the prompt as the high-level prior and the
diffusion model as the perceptual substrate, and lowering g weakens a top-down
prior, the direction REBUS assigns to psychedelics
\citep{carhartharris2019rebus}. Read the diffusion prior as the prior, and
lowering g is prior dominance, closer to sensory-deprivation and Charles Bonnet
accounts. The data bear on that choice in one way only: the empty-conditioning
baseline is the extreme of breakdown on both checkpoints, and on SDXL an empty
prompt is still an embedding, not the zero-conditioning branch. The neuroscience
is the source of the rubric and the framing. Nothing below tests the analogy.

The assay is five lines. I generated a fixed grid, 6 prompts by 7 guidance
values by 10 seeds on two checkpoints, 860 images with empty-prompt baselines.
Two judges from different model families scored every image blind on the
Klüver-derived rubric, guidance value hidden and order shuffled. Two automatic
``does this look like a good image'' scores per image were pre-registered as
quality proxies, so that the breakdown association could be re-estimated after
adjustment for those scores. Metric, direction, disconfirming outcome and four
gates were committed to git before any analysed run; the whole tree is public.
Under a two-judge panel amended after the first run, lower classifier-free
guidance was associated with greater rubric-scored visual breakdown, and the
association remained after adjustment for two no-reference quality proxies. SDXL
met every pre-registered gate. SD 3.5 missed the slope gate: its slope of -0.182
sits above the -0.20 line, and its interval straddles that line, so an effect of
the registered size is neither shown nor ruled out. SDXL's slope is steeper, and
the difference was tested. Looking afterwards, almost all of SD 3.5's drop
happens between g = 1 and g = 2.

Four contributions.

\begin{enumerate}
\def\labelenumi{\arabic{enumi}.}
\tightlist
\item
  \textbf{The assay and its pre-registration record.} One knob, a fixed grid, a
  blind cross-family judge panel, a pre-specified quality adjustment, and
  confirming and disconfirming criteria in git before the run, with one dated
  pre-data amendment and one post-data instrument amendment. A git log verifies
  the order.
\item
  \textbf{The dose-response, and what separates the two checkpoints tested.} The
  slope, whether the four scores fall in a straight line as guidance rises, in
  standardised units, is -0.340 on SDXL and -0.182 on SD 3.5, against a
  registered gate of -0.20. The slope difference was tested with a model by
  guidance interaction, not read off the pair. On the exploratory veridicality
  contrast, 64\% of the distortion field's absolute correlation with guidance
  attenuated after adjustment for veridicality; that is construct overlap, not a
  causal split.
\item
  \textbf{Silent judge failure, and the screen that catches it.} Two
  vision-language judges returned zero on every field of every probe image while
  emitting well-formed JSON and fluent, confident captions. The panel change
  that followed is a post-data amendment, not a pre-specified drop-in. Refilling
  one seat with no rubric change raised composite inter-judge kappa, agreement
  between two judges beyond chance where 0 is chance and 1 is identical scoring,
  on SDXL from 0.290 to 0.562; Section 3.4 reports both checkpoints.
\item
  \textbf{The form-constant result.} The registered validation of the Exp 01
  judge failed on specificity: its class flags fire on 45\% of ordinary
  generated images, above the ceiling of 20\% fixed in advance. Under the rule
  written into that pre-registration, the Exp 01 null is uninterpretable, and I
  say so. The sensitivity reading is exploratory, and so is the separation
  between the rendered stimuli and ordinary scenes on the geometric-intensity
  metric Exp 01 used.
\end{enumerate}

\subsection{Related work}\label{related-work}

Ho and Salimans introduced classifier-free guidance and its trade-off: more
guidance raises sample fidelity and reduces diversity \citep{ho2022cfg}. Saharia
et al.~traced oversaturated, unnatural images at high guidance weights to a
train-test mismatch, and introduced dynamic thresholding to keep large weights
usable \citep{saharia2022imagen}. Kynkäänniemi et al.~showed that the effect of
guidance depends on the noise level, harmful at high noise and largely
unnecessary at low noise, and described low guidance as yielding diverse but
fuzzy images that lack detail \citep{kynkaanniemi2024interval}. Karras et
al.~described guidance as pulling samples toward well-learned, high-probability
regions, and guided a model with an inferior version of itself instead
\citep{karras2024autoguidance}. Suzuki et al.~built the Hallucination Machine,
applying DeepDream to panoramic video and comparing participants' ratings
against psilocybin questionnaire data \citep{suzuki2017}. Suzuki set out
computational neurophenomenology and three axes, veridicality, spontaneity and
complexity, as phenomenology first, not a direct map to pathological mechanism
\citep{suzuki2024}. Chen et al.~found that multimodal judges track human
preference on pair comparison while diverging from it on scoring evaluation,
that the other models they evaluate (LLaVA, CogVLM, Qwen-VL-Max) align worse
with human judgement than GPT-4V, and that judge outputs carry biases and
hallucinated justifications \citep{chen2024mllmjudge}. Xie et al.~showed that
common human-preference evaluators are biased toward large guidance scales, so
raising the guidance scale improves preference scores even where image quality
is damaged, and proposed a guidance-aware evaluation in response
\citep{xie2026guidancematters}.

Lowering guidance is not the new thing here. Kynkäänniemi et al.~study low and
zero guidance already. What is new is the rubric that names failure phenomena
rather than rating preference or quality, the pre-registration of the metric and
the gates, the comparison across two checkpoints, and the judge audit. Xie et
al.~score preference and repair the evaluator for its guidance bias; I score
named phenomena instead of preference, so that work is a complement here and not
a competitor.

\section{The assay}\label{the-assay}

\subsection{Overview: two guidance instantiations and a calibration
arm}\label{overview-two-guidance-instantiations-and-a-calibration-arm}

The contribution is a \textbf{measurement procedure}: take a generative system
that works, loosen one control parameter, classifier-free guidance, hold
everything else fixed, and measure what breaks. The guidance assay is
instantiated twice, in Exp 01 on a single prompt and in Exp 03 on six prompts.
Exp 02 is not a guidance sweep at all: it is a calibration arm, instrument
validation on a neural field, and what it validates is the judge that scored Exp
01.

Four steps, identical across the two guidance instantiations:

\begin{enumerate}
\def\labelenumi{\arabic{enumi}.}
\tightlist
\item
  \textbf{Generate} a full grid (guidance x prompt x seed) with all else pinned,
  plus an empty-prompt baseline at the same seeds.
\item
  \textbf{Score blind.} A judge sees pixels and a fixed rubric naming the
  intended scene. The guidance value is never shown, order is shuffled, decoding
  is deterministic, and un-blinding happens only at analysis.
\item
  \textbf{Analyse against a plan committed before the run}: metric, direction,
  confirming and disconfirming outcomes, and test.
\item
  \textbf{Publish the whole tree}: pre-registration, judgements, raw judge
  replies, reports, images.
\end{enumerate}

\subsection{Pre-registration as git
history}\label{pre-registration-as-git-history}

Each experiment's hypothesis, sweep, held-fixed variables, metric and direction,
the disconfirming null, sample size and test are committed \textbf{before any
run on the configuration to be analyzed}. A later change is a dated commit
reclassifying the experiment as exploratory:

\begin{verbatim}
git log -- experiments/03_l23_hardening/preregistration.json
\end{verbatim}

Exp 03 carries two dated amendments. The first (2026-07-21) is pre-data: Judge A
moved to the Batches API, a further judge family was added, and reliability was
generalized to the mean of pairwise weighted kappas, leaving rubric, prompts,
grid, endpoint and thresholds untouched. Being pre-data is its whole defence:
when it was committed only a two-image smoke test had been run, so no image
entering any analysis below existed. The second (2026-08-08) is post-data, and
it is an instrument amendment: after the full run of 2026-07-22 had been scored,
the judge B seat was refilled, under the pre-registration clause listing judge
models as swappable without reclassifying the experiment. It leaves rubric text,
prompts, grid, endpoint and thresholds untouched, the trigger condition the
pre-registration had named was met, and the model that took the seat is a
deviation from the fallback named there. Section 3.4 reports what forced it and
what it changed.

\subsection{Models and generation}\label{models-and-generation}

\begin{longtable}[]{@{}
  >{\raggedright\arraybackslash}p{(\linewidth - 4\tabcolsep) * \real{0.3333}}
  >{\raggedright\arraybackslash}p{(\linewidth - 4\tabcolsep) * \real{0.3333}}
  >{\raggedright\arraybackslash}p{(\linewidth - 4\tabcolsep) * \real{0.3333}}@{}}
\caption{Generation design of the two guidance instantiations. Everything not
listed is the pipeline default and is pinned in each experiment's
metadata.}\tabularnewline
\toprule\noalign{}
\begin{minipage}[b]{\linewidth}\raggedright
\end{minipage} & \begin{minipage}[b]{\linewidth}\raggedright
Exp 01
\end{minipage} & \begin{minipage}[b]{\linewidth}\raggedright
Exp 03
\end{minipage} \\
\midrule\noalign{}
\endfirsthead
\toprule\noalign{}
\begin{minipage}[b]{\linewidth}\raggedright
\end{minipage} & \begin{minipage}[b]{\linewidth}\raggedright
Exp 01
\end{minipage} & \begin{minipage}[b]{\linewidth}\raggedright
Exp 03
\end{minipage} \\
\midrule\noalign{}
\endhead
\bottomrule\noalign{}
\endlastfoot
Architectures & SDXL base 1.0 (UNet) \citep{podell2023sdxl}, SD 3.5 medium
(MMDiT) \citep{esser2024sd3} & same two \\*
Guidance grid & 9 values: 1.0, 1.5, 2.0, 3.0, 4.5, 6.0, 8.0, 11.0, 15.0 & 7
values: 1.0, 2.0, 3.0, 5.0, 7.0, 11.0, 15.0 \\*
Prompts & \textbf{1} (the still life; becomes \texttt{p1\_stilllife} in Exp 03)
& 6, including a designed control \\*
Seeds & 10 (42-51) & 10 (42-51), held fixed across guidance and model \\*
Baseline & empty prompt, same seeds & empty prompt, same seeds \\*
Steps / size & 25 / 1024x1024 & 25 / 1024x1024 \\*
Total images & 100 per model, 200 total & 430 per model, 860 total \\*
\end{longtable}

A note on wording used throughout: I write \emph{checkpoint}, not
\emph{architecture}. A checkpoint is one trained model, a single file of learned
weights. SDXL and SD 3.5 do use different architectures (UNet and MMDiT), but
each architecture appears here as exactly one checkpoint, and that checkpoint
also differs from the other in training data, training objective, text encoder
and sampler. Any difference between the two therefore belongs to these two
trained models; it cannot be credited to the architecture without at least a
second checkpoint of each.

Held fixed on Apple Silicon: the model-default scheduler, \texttt{force\_upcast}
at the SDXL VAE decode, and float32 for SDXL, which in float16 on MPS gives
all-NaN latents above g = 1.

Conventions differ by one, and the difference matters for reading the grid.
These experiments use the diffusers convention, where the sampler forms
\texttt{pred\ =\ uncond\ +\ g\ (cond\ -\ uncond)}, so g = 1 is the unguided
conditional model; Ho and Salimans write the same combination with a weight w
for which w = 0 is unguided \citep{ho2022cfg}. Every guidance value here is g.

Exp 01 is a \textbf{single-prompt} design, so a guidance effect and a property
of that scene are inseparable; Exp 03's six prompts are what make the per-prompt
breakdown possible. They are fixed in the pre-registration, with their ids as
used throughout:

\begin{itemize}
\tightlist
\item
  \texttt{p1\_stilllife}, a watermelon, a glass half-filled with water, and a
  set of keys on a wooden table (the Exp 01 prompt verbatim, kept for
  continuity)
\item
  \texttt{p2\_portrait}, a close-up portrait photograph of an elderly man's face
\item
  \texttt{p3\_bicycle}, a red bicycle leaning against a brick wall
\item
  \texttt{p4\_oranges}, a bowl of oranges on a kitchen counter
\item
  \texttt{p5\_livingroom}, an empty living room with a sofa, a lamp, and a
  window
\item
  \texttt{p6\_forest}, a forest of pine trees on a misty mountainside
\end{itemize}

The forest prompt is the designed control, pre-registered as the low-objecthood
arm: a scene with no small fixed inventory of distinct foreground objects, where
reduplication has little to multiply and tiling was the predicted stress. Its
pre-registered role reads ``low-objecthood control; should show tiling but LOW
reduplication/condensation if the effect is object-bound'', so it is the prompt
on which an object-bound account of the breakdown can come apart.

\subsection{The two rubrics}\label{the-two-rubrics}

Four phenomena from Klüver's inventory of hallucinatory constants
\citep[pp.~175-207]{kluver1966} are scored per image on a 0-3 ordinal scale:
\textbf{reduplication} (his polyopia, one form multiplied), \textbf{distortion}
(his dysmorphopsia, a form melted or made impossible), and
\textbf{fragmentation} and \textbf{condensation}, two of his spatio-temporal
transformations. A binary \textbf{tiling} flag records dissolution into a
repeating field. Dysmegalopsia, the size change completing his triad, is
dropped. The rubric names the intended scene and never the guidance value.

The Klüver fields describe what happens to objects locally: copies, pieces,
fusions, warps. A second rubric describes the image as a whole. Suzuki,
Schwartzman and Seth \citep{suzuki2024} compared visual hallucinations with
different causes (neurodegenerative, visual loss, psychedelic) and found they
differ along three continuous axes: \textbf{veridicality}, how real and coherent
the hallucinated content looks; \textbf{spontaneity}, how independent it is of
what is actually in front of the eyes; and \textbf{complexity}, from simple
geometry to elaborate scenes. They reproduced those differences in deep networks
by varying which network layer drives the image and whether a learned image
prior constrains it, a manipulation close in spirit to turning guidance. Exp 03b
scores the same images on those three axes, 0 to 3 each, with a binary
coherent-scene flag, carried over from Exp 01. It asks whether a global reading
moves with guidance the way the local one does, and whether it carries signal
the Klüver fields do not. Both judge prompts are given verbatim in the appendix.
The adaptation changes two of the three: spontaneity here is content the prompt
did not ask for rather than independence from sensory input, and veridicality
here combines realism and coherence rather than realism alone. Its standing is
weaker: the Klüver scores were known when it was committed, so its predictions
are pre-data with respect to the axes scores only.

\subsection{Judges}\label{judges}

The panel reported throughout is two model judges from different families, plus
two human raters on a blind subset, the author and a second rater who does not
know the hypothesis. It is not the panel that ran first. The full run of
2026-07-22 used Claude Sonnet 5, Qwen2.5-VL-7B in the judge B seat and
Llama-3.2-11B-Vision in a judge C seat. The human subset and the screening
probes described below were run after that run, on 2026-08-08, and exposed two
inert raters; the judge B seat was then refilled with Qwen3-VL-32B and judge C
archived. The pre-registration had named Qwen2.5-VL-32B as the fallback if the
7B's human-subset kappa were poor. That trigger condition was met, but the model
actually used is Qwen3-VL-32B, one generation newer, because it was the
candidate that cleared the screen and the named model was not usable; the
amendment records the substitution as a deviation from the named fallback. Which
candidates survived the screen is a result, reported in Section 3.4.

\begin{longtable}[]{@{}
  >{\raggedright\arraybackslash}p{(\linewidth - 2\tabcolsep) * \real{0.5000}}
  >{\raggedright\arraybackslash}p{(\linewidth - 2\tabcolsep) * \real{0.5000}}@{}}
\caption{Judge panel as reported, the pre-registered original panel, and the
candidates screened.}\tabularnewline
\toprule\noalign{}
\begin{minipage}[b]{\linewidth}\raggedright
\end{minipage} & \begin{minipage}[b]{\linewidth}\raggedright
judge
\end{minipage} \\
\midrule\noalign{}
\endfirsthead
\toprule\noalign{}
\begin{minipage}[b]{\linewidth}\raggedright
\end{minipage} & \begin{minipage}[b]{\linewidth}\raggedright
judge
\end{minipage} \\
\midrule\noalign{}
\endhead
\bottomrule\noalign{}
\endlastfoot
Judge A & Claude Sonnet 5, via the Message Batches API \\*
Judge B & Qwen3-VL-32B local (MLX), which replaced the pre-registered
Qwen2.5-VL-7B on 2026-08-08; a deviation from the named Qwen2.5-VL-32B
fallback \\*
screened, not used & Qwen2.5-VL-7B, Qwen3-VL-8B, Llama-3.2-11B-Vision,
Gemma-3-27B \\*
Human & two raters (author; naive second rater), 28 images \\*
\end{longtable}

All judges receive the identical rubric from one module under identical
blinding, shuffling and decoding; raw replies are retained beside the parsed
record.

\textbf{Judge screening (\texttt{RESEARCH\_METHODOLOGY.md} section 5.3, in the
repository).} The rule is that a judge is screened per rubric \emph{before} the
full run, on a stratified subset of 28 images spanning architecture, guidance
bin and prompt, under the blinding and decoding of the full pass. The rule is
what this experiment produced rather than what it followed: it was added to the
methodology on 2026-08-09, and in Exp 03 itself the Klüver-rubric screen was run
post-data, on 2026-08-08, after the full pass of 2026-07-22 had been scored. The
probe reports gradedness, the share of scores landing on an interior point of
the scale rather than an endpoint, and flags any graded field whose score never
varies; a judge flat on a graded field is disqualified for that rubric. The
screen runs on the four 0-3 fields only; the binary tiling flag is excluded from
it. The panel was probed before the axes pass, which is the rule applied as
written. Selection is on gradedness alone; human agreement is reported
afterwards, not used to choose.

\subsection{Instrument validation (Exp 02)}\label{instrument-validation-exp-02}

Exp 02 answers an objection to Exp 01: \emph{could the judge detect a form
constant at all?} A minimal scalar neural field in the Ermentrout-Cowan family
\citep{ermentrout1979}, rendered through the inverse log-polar map, is driven
through the Turing bifurcation with its kernel and that map held fixed and the
critical point derived by linear stability analysis rather than tuned.
Bressloff's orientation-field extension \citep{bressloff2001} was not run. Its
renderings go through the detector, rubric text and decoding used on the
generated images, mixed with blank fields and ordinary scenes, so only the image
changes between control and experiment. It was pre-registered as confirmed only
if all four judge class labels are realized, blanks stay below baseline, and the
confusion matrix is diagonal-dominant. The same registration fixed the test of
the judge itself: it had to flag at least 80\% of the rendered form constants
and at most 20\% of ordinary generated images and of blank fields, and if it
failed, the Exp 01 null was to be reported as uninterpretable.

\subsection{Endpoints and statistics}\label{endpoints-and-statistics}

\textbf{Primary endpoint (Exp 03).} The endpoint is the outcome measure the
hypothesis is tested on. Per image and field, the judge-mean 0-3 score, z-scored
across all conditioned images \emph{within a model}; the composite is the
per-image mean of the four z-scored fields. The endpoint is therefore
judge-scored breakdown on the Klüver-derived rubric, not human-perceived
breakdown: its evidence of human validity is two raters on 28 images, who agree
with each other at composite kappa 0.567 and with the better judge at 0.337 and
0.426 (Section 3.4).

\textbf{Primary model.} The 420 images per checkpoint (SDXL or SD 3.5) are not
independent: seventy share each prompt, and a busy forest scene scores more
breakdown than a portrait at every guidance value. A plain regression would let
those prompt differences leak into the guidance effect. A linear mixed model
separates them. It fits one slope for guidance, shared by all images, and lets
each prompt and each seed have its own baseline level of breakdown:

\begin{verbatim}
composite ~ guidance_std + (1 | prompt) + (1 | seed)
\end{verbatim}

The slope therefore measures how breakdown changes with guidance within the same
prompt and seed, not between them. Technically, guidance is z-scored and treated
as interval, the slope is the fixed effect, and the per-prompt and per-seed
baselines are random intercepts.

\textbf{Quality adjustment.} Breakdown at low guidance could be low-quality
artifact rather than objecthood dissolution, so a no-reference quality score is
computed per image: CLIP-IQA \citep{wang2023clipiqa} and a LAION aesthetic
predictor are each z-scored across the model's conditioned images and averaged
into one score, Q, on the same standardised ruler as the composite. The guidance
association is then re-estimated as a partial correlation adjusting for Q. The
pre-registration named a quality-matched two-arm contrast; the implementation
trims both arms to the overlapping quality band rather than matching images pair
by pair or equalising mean quality. That is an implementation deviation,
labelled as common-support trimming wherever it is reported. Quality is measured
on the same images after guidance has acted, so the partial correlation is a
sensitivity analysis, not a causal de-confound. An adjustment chosen afterwards
cannot be told from a defence; this one was pre-registered.

\textbf{Multiplicity.} The primary family is two tests, one per architecture;
the secondary family (4 fields x 2 architectures) is Benjamini-Hochberg
corrected \citep{benjamini1995fdr}, which controls the expected share of false
positives among the tests that pass, and is reported as q-values, the corrected
p-values, with effect sizes.

\textbf{Reliability.} Quadratic-weighted Cohen's kappa \citep{cohen1968kappa},
percent agreement, and Gwet's AC2 \citep{gwet2008ac2}, an agreement coefficient
that stays stable when most scores sit at zero, where kappa is pulled down; each
is averaged over judge pairs. The composite kappa that the gate is written on is
the mean of the four per-field quadratic-weighted kappas, over the four
intensity fields with tiling excluded; it is called the composite (mean
per-field) kappa below. The pre-registration does not let kappa be read alone,
so AC2 and percent agreement are reported beside it.

\textbf{Confirmation, stated in advance.} Every gate below had to clear \emph{on
both architectures}; the design can miss at any one. In words, the four gates
are four different questions.

The first, the slope, is whether there is an effect at all: does breakdown fall
as guidance rises, by at least the registered amount, with an interval that does
not touch zero.

The second, the partial correlation, is the quality adjustment. The worry is
this: when guidance is low, the images are also lower quality in a general
sense, blurrier, less polished. A judge looking at a blurry image might call it
fragmented or distorted just because it is blurry, not because objects actually
broke apart. If that were the whole story, the paper would have measured nothing
but ``low guidance makes worse pictures'', which everyone already knew. So every
image also gets a quality score from two separate automatic raters that know
nothing about the rubric; they only answer ``does this look like a good image''.
The second gate then asks: after adjustment for those two proxies, does
breakdown still go down as guidance goes up? The number that answers it is the
partial correlation. If breakdown were nothing more than what those two proxies
measure, that number would drop to zero. The proxies are not a matched quality
design: they do not hold image quality fixed, and they cannot rule out kinds of
degradation they do not measure.

The third, prompt generality, is whether it is more than one prompt: at least
five of the six prompts have to show the effect on their own.

The fourth, kappa, is not about guidance at all. It only asks whether the two
judges gave similar scores to the same images. If one judge says distortion 3
and the other says distortion 0 on the same picture, the score is noise and
nothing built on it means anything. Kappa is agreement between the two judges
beyond what chance would give. It is a check on the instrument, done before
anything the instrument reads is trusted.

So the four gates are: is there an effect, is it more than those two quality
proxies, is it more than one prompt, and can the measuring device be trusted at
all.

\begin{longtable}[]{@{}ll@{}}
\caption{Confirmation criteria committed before the analysed run. All four had
to clear on both checkpoints.}\tabularnewline
\toprule\noalign{}
gate & pre-registered threshold \\
\midrule\noalign{}
\endfirsthead
\toprule\noalign{}
gate & pre-registered threshold \\
\midrule\noalign{}
\endhead
\bottomrule\noalign{}
\endlastfoot
standardised composite slope on guidance & \textless= -0.20, CI excluding
zero \\*
partial rho adjusting for two quality proxies & \textless= -0.20, CI excluding
zero \\*
prompts with a negative per-prompt slope & \textgreater= 5 of 6 \\*
inter-judge composite weighted kappa & \textgreater= 0.4 \\*
architectures required to clear all four & both \\*
\end{longtable}

Three rules, dated in \texttt{RESEARCH\_METHODOLOGY.md}, constrain the analyses:
shape-agnostic endpoints, no pooled correlation without its per-prompt
breakdown, and a judge screened per rubric. What they cost is revisited in
Section 4.

\subsection{Reproducibility}\label{reproducibility}

Code, pre-registrations, analyses, judge outputs and images are public, and
every figure and number below regenerates from them with one command per
experiment. Section 6 gives the repository, the datasets and the commands.

\begin{center}\rule{0.5\linewidth}{0.5pt}\end{center}

\section{Results}\label{results}

\subsection{Dose-response on the
composite}\label{dose-response-on-the-composite}

\begin{figure}
\centering
\includegraphics[width=0.92\linewidth,height=\textheight,keepaspectratio,alt={The phenomenon, without choosing the best case. The living-room prompt on SDXL (top) and SD 3.5 (bottom). Left: guidance g = 1, the sixth of the ten seeds ranked by judge-scored composite, one of the two middle ranks since ten seeds have no single median (SDXL seed 42, SD 3.5 seed 48). Middle: g = 1, the seed the judges score highest (SDXL seed 44, SD 3.5 seed 50), the clearest case rather than a typical one. Right: g = 11, the sixth-ranked seed at that level (SDXL seed 50, SD 3.5 seed 51). Within a row the prompt and the steps are identical. On SDXL the middle-ranked low-guidance image is already a sketch with fractured furniture; on SD 3.5 it is a plausible room with melted objects and a style shift, which is less than the SDXL row shows and is why the two checkpoints are compared by slope and not by eye.}]{figures/fig1_livingroom_pair.png}
\caption{The phenomenon, without choosing the best case. The living-room prompt
on SDXL (top) and SD 3.5 (bottom). Left: guidance g = 1, the sixth of the ten
seeds ranked by judge-scored composite, one of the two middle ranks since ten
seeds have no single median (SDXL seed 42, SD 3.5 seed 48). Middle: g = 1, the
seed the judges score highest (SDXL seed 44, SD 3.5 seed 50), the clearest case
rather than a typical one. Right: g = 11, the sixth-ranked seed at that level
(SDXL seed 50, SD 3.5 seed 51). Within a row the prompt and the steps are
identical. On SDXL the middle-ranked low-guidance image is already a sketch with
fractured furniture; on SD 3.5 it is a plausible room with melted objects and a
style shift, which is less than the SDXL row shows and is why the two
checkpoints are compared by slope and not by eye.}\label{fig:phenomenon}
\end{figure}

Of the 860 images generated, 418 (SDXL) and 417 (SD 3.5) conditioned images
entered the primary analysis, after listwise removal of every image with a
failed judge reply, 2 and 3 Claude parse failures respectively, together with 10
empty-prompt baselines per model.

Table 4 reports the four pre-registered gates. SDXL clears all four, its six
per-prompt Spearman correlations of composite against guidance all negative,
from -0.74 (bicycle) to -0.28 (portrait). SD 3.5 misses the slope gate:
per-prompt correlations from -0.49 (living room) to +0.30 (portrait), and a
slope whose point estimate misses the registered gate of -0.20 while its
interval includes it. Its kappa of 0.440 is computed on the registered full
complete-case set, including ten empty-prompt baselines (n = 427); on
conditioned images only it is 0.394, below the 0.4 gate. I treat the registered
full-set number as the gate and report the conditioned-only value as a
sensitivity. The joint claim, requiring every gate on both architectures, is not
confirmed. The forest prompt, the pre-registered low-objecthood control, is
counted inside both of those tallies; it is read as a control below rather than
as a sixth ordinary prompt.

\begin{longtable}[]{@{}lllll@{}}
\caption{Pre-registered gates, amended panel (Claude Sonnet 5 and Qwen3-VL-32B).
Slope is the guidance-standardised fixed effect of guidance on the four-field
composite from a linear mixed model with random prompt and seed intercepts, 95\%
CI. Partial rho adjusts for the two-proxy no-reference quality score. Prompts is
the count with a negative per-prompt Spearman. Kappa is the quadratic-weighted
composite between judges on all complete images, including ten empty-prompt
baselines per checkpoint.}\tabularnewline
\toprule\noalign{}
& slope (95\% CI) & partial rho (95\% CI) & prompts negative & kappa \\
\midrule\noalign{}
\endfirsthead
\toprule\noalign{}
& slope (95\% CI) & partial rho (95\% CI) & prompts negative & kappa \\
\midrule\noalign{}
\endhead
\bottomrule\noalign{}
\endlastfoot
pre-registered threshold & \(\le\) -0.20 & \(\le\) -0.20 & \(\ge\) 5/6 & \(\ge\)
0.4 \\*
SDXL & -0.340 {[}-0.400, -0.279{]} & -0.433 {[}-0.508, -0.347{]} & 6/6 &
0.562 \\*
SD 3.5 & -0.182 {[}-0.248, -0.115{]} & -0.245 {[}-0.338, -0.146{]} & 5/6 &
0.440 \\*
\end{longtable}

The pre-registered secondary family, four fields by two architectures with
Benjamini-Hochberg correction, moves the same way: on SDXL every field falls
with guidance, fragmentation -0.618, distortion -0.486, condensation -0.442 and
reduplication -0.116, all at q \textless= 0.018, and on SD 3.5 every field falls
as well, from fragmentation -0.453 to distortion -0.114, all four clearing q
\textless= 0.05.

Kappa is not read alone. Across the same four fields on the amended panel,
Gwet's AC2 runs from 0.710 to 0.924 on SDXL and from 0.803 to 0.953 on SD 3.5,
and exact percent agreement from 38\% to 72\% and from 38\% to 82\%.
Fragmentation is the top of every one of those ranges, and distortion the bottom
of three of the four.

After adjustment for the two quality proxies the association remains on both
checkpoints (Table 4). On SDXL it also holds when the analysis is repeated on
resampled sets of prompts (a prompt-level bootstrap, Section 5); on SD 3.5 it
does not. Without the adjustment the correlations are -0.463 on SDXL and -0.216
on SD 3.5, so it barely moves either: had breakdown been nothing more than what
the two proxies measure, the partial correlation would have fallen toward zero.
Outside the mixed model, on the common-support-trimmed arms, Cliff's delta is
0.49 {[}0.362, 0.606{]} on SDXL and 0.188 {[}0.050, 0.320{]} on SD 3.5. Trimming
to the overlapping quality band does not equalise the arms: the standardised
mean difference in quality between the high- and low-guidance arms that survives
the trim is 0.740 on SDXL and 0.284 on SD 3.5, so these are common-support arms
and not matched ones, a deviation from the registered matching language. A
prompt-cluster bootstrap of the quality-adjusted partial correlation, reported
in Section 5, is the sharper test of how far the adjustment carries.

Refit on each judge separately, both slopes stay negative: Claude -0.321 and
Qwen -0.295 on SDXL, Claude -0.174 and Qwen -0.141 on SD 3.5.

A note on fitting. The mixed model has five numbers to estimate: the guidance
slope, an overall baseline, how much baselines vary between prompts, how much
they vary between seeds, and the remaining noise. They are found by numerical
search, here with the \texttt{MixedLM} routine of the Python statsmodels
library, which reports whether the search converged, that is, confirmed it had
reached the best values. On SD 3.5 the default search (L-BFGS) stopped without
converging; the best value for the seed-to-seed variation is zero, the edge of
what a variance can be, a common difficulty for that method. A second method
(Powell) converged and gives the same slope to three decimals, -0.182 {[}-0.248,
-0.115{]}; that converged fit is the one reported. On SDXL the default search
converged.

The designed control is read post-hoc and at the field level (Exp 03
\texttt{analysis.md} section 9.2, in the repository), because the
pre-registration gated on the composite only. On SDXL the forest composite still
tracks guidance, rho -0.62, so the control did not stay flat and the
object-bound reading of the breakdown is not established by it. What did behave
as pre-registered is the pair of fields the control was built to dissociate:
reduplication is flat at +0.05, and tiling at the lowest guidance value reaches
40\%, the highest rate of any prompt. On SD 3.5 the forest composite is flat at
-0.02, and that flatness is cancellation between fields rather than absence of
movement: reduplication runs +0.52 against fragmentation -0.38. The control
therefore separates the fields it was designed to separate without vindicating
the composite-level prediction.

Two features precede any modelling of shape (Figure 2). The empty-prompt
baseline is the extreme of the composite on both architectures, above even g =
1: 1.51 standardised units on SDXL and 2.17 on SD 3.5. And the composite is not
the quality score renamed: the quality curves differ in shape between the two
checkpoints while both composites fall.

\begin{figure}
\centering
\pandocbounded{\includegraphics[keepaspectratio,alt={Dose-response. The four-field judge-scored rubric breakdown composite (circles) and the no-reference quality score (squares), in standardised units within each checkpoint, against guidance on a log axis. The shaded band is a bootstrap 95\% interval of the mean over images at each guidance value; the thin lines are the six per-prompt means, which is the spread every pooled number in the text travels with. The empty-prompt arm, having no guidance value, is detached at the right. Each point is a mean over complete cases at that guidance value, 58 to 60 images after parser failures.}]{figures/fig2_dose_response.png}}
\caption{Dose-response. The four-field judge-scored rubric breakdown composite
(circles) and the no-reference quality score (squares), in standardised units
within each checkpoint, against guidance on a log axis. The shaded band is a
bootstrap 95\% interval of the mean over images at each guidance value; the thin
lines are the six per-prompt means, which is the spread every pooled number in
the text travels with. The empty-prompt arm, having no guidance value, is
detached at the right. Each point is a mean over complete cases at that guidance
value, 58 to 60 images after parser failures.}\label{fig:dose}
\end{figure}

\subsection{Shape across the two
checkpoints}\label{shape-across-the-two-checkpoints}

Table 4 splits the two checkpoints: SDXL clears the slope gate and SD 3.5 misses
it narrowly. This section asks whether that split is real and, if so, what it
consists of. Figure 2 shows both curves; Figure 3 puts the same averages on a
plain guidance axis and draws the registered straight-line slope through them.

\begin{figure}
\centering
\includegraphics[width=0.92\linewidth,height=\textheight,keepaspectratio,alt={Why a straight line understates SD 3.5. Mean breakdown composite at each guidance value (points, about 60 images each) on a linear guidance axis, with the registered straight-line slope (solid) and the slope refitted on log2 guidance (dashed; a sensitivity analysis pre-specified with no gate). Both lines are drawn from the fitted slopes through the overall mean. On SDXL the straight line is a fair summary; on SD 3.5 almost all of the fall is between g = 1 and g = 2, which a straight line cannot follow.}]{figures/fig3_slopes.png}
\caption{Why a straight line understates SD 3.5. Mean breakdown composite at
each guidance value (points, about 60 images each) on a linear guidance axis,
with the registered straight-line slope (solid) and the slope refitted on log2
guidance (dashed; a sensitivity analysis pre-specified with no gate). Both lines
are drawn from the fitted slopes through the overall mean. On SDXL the straight
line is a fair summary; on SD 3.5 almost all of the fall is between g = 1 and g
= 2, which a straight line cannot follow.}\label{fig:slopes}
\end{figure}

On SDXL the straight line is a fair summary: breakdown keeps falling as guidance
rises. On SD 3.5 it is not: almost all of the fall happens between g = 1 and g =
2, and the curve is flat after that. A straight line drawn through a cliff tilts
less than the cliff, so the registered linear slope understates an effect packed
into one step. Refitted against log2 of guidance, which follows a cliff more
closely, the SD 3.5 slope is -0.301 instead of -0.182 (Table 5). That refit is a
sensitivity analysis, pre-specified with no gate after the verdict was fixed,
and it does not change the verdict. In hindsight a linear endpoint on this grid
was the wrong pre-registration: it treats the step from g = 1 to g = 2, where
most of the change happens, as one fourteenth of the range.

So SD 3.5's miss is narrow, and partly a matter of shape. Its slope is -0.182
{[}-0.248, -0.115{]}. The interval excludes zero, so the effect is real; the
point estimate misses the -0.20 gate and the interval includes it, so an effect
of the registered size is neither shown nor ruled out. That is not a null: the
word is reserved here for the Exp 01 form-constant result of Section 3.5, which
sits inside its pre-registered null region.

Is the split real, or did two noisy estimates of the same effect simply land on
either side of a line? Comparing -0.340 and -0.182 by eye ignores the
uncertainty in each, so the difference was tested directly: one mixed model
fitted on both checkpoints together, with a term that lets the guidance slope
differ between them (a model by guidance interaction, pre-specified on
2026-09-17 and exploratory). SDXL's slope is steeper by -0.158 {[}-0.248,
-0.067{]} on the pre-registered z-composite and by -0.164 {[}-0.222, -0.105{]}
on the raw judge-mean scores, so the z-scoring did not create the difference.
The two checkpoints respond differently to guidance. Each architecture is
represented by one checkpoint, which also differs in training data, objective,
text encoder and sampler, so the difference belongs to the two checkpoints
tested and is not shown to be architectural.

\begin{longtable}[]{@{}
  >{\raggedright\arraybackslash}p{(\linewidth - 4\tabcolsep) * \real{0.3333}}
  >{\raggedright\arraybackslash}p{(\linewidth - 4\tabcolsep) * \real{0.3333}}
  >{\raggedright\arraybackslash}p{(\linewidth - 4\tabcolsep) * \real{0.3333}}@{}}
\caption{Sensitivity and post-hoc analyses of the guidance effect. Row 1 is the
pre-registered primary slope. Rows 2 to 5 are mixed-model sensitivities
pre-specified on 2026-09-03 with no gate. The last two rows are post-hoc shape
descriptors from posthoc.py.}\tabularnewline
\toprule\noalign{}
\begin{minipage}[b]{\linewidth}\raggedright
\end{minipage} & \begin{minipage}[b]{\linewidth}\raggedright
SDXL
\end{minipage} & \begin{minipage}[b]{\linewidth}\raggedright
SD 3.5
\end{minipage} \\
\midrule\noalign{}
\endfirsthead
\toprule\noalign{}
\begin{minipage}[b]{\linewidth}\raggedright
\end{minipage} & \begin{minipage}[b]{\linewidth}\raggedright
SDXL
\end{minipage} & \begin{minipage}[b]{\linewidth}\raggedright
SD 3.5
\end{minipage} \\
\midrule\noalign{}
\endhead
\bottomrule\noalign{}
\endlastfoot
linear g (pre-registered) & -0.340 {[}-0.400, -0.279{]} & -0.182 {[}-0.248,
-0.115{]} \\*
log2 g & -0.440 {[}-0.494, -0.386{]} & -0.301 {[}-0.362, -0.239{]} \\*
rank of g & -0.415 {[}-0.471, -0.359{]} & -0.265 {[}-0.328, -0.202{]} \\*
three fields, no distortion, linear g & -0.308 {[}-0.365, -0.251{]} & -0.206
{[}-0.269, -0.144{]} \\*
three fields, no distortion, log2 g & -0.395 {[}-0.447, -0.344{]} & -0.307
{[}-0.365, -0.248{]} \\*
drop g = 1, Spearman composite vs g (post-hoc) & rho -0.244, p \textless{} 0.001
& rho 0.048, p = 0.3757 \\*
share of the g = 1-to-curve-minimum drop occurring from g = 1 to 2 (post-hoc) &
60\% & 91\% \\*
\end{longtable}

Two post-hoc numbers in Table 5 put the shape difference in figures. The first
drops g = 1: the correlation is recomputed without the g = 1 images, which asks
whether any effect remains above the bottom step of the dial. On SDXL it does;
on SD 3.5 it does not. And of each curve's drop from g = 1 to its lowest point,
60\% happens in the first step on SDXL against 91\% on SD 3.5. Both are
post-hoc, and g = 1 is the dose level hardest to separate from
under-conditioning artifact. The other sensitivity rows agree in direction;
under them SD 3.5 is negative on six of six prompts, its portrait prompt at
-0.42 against +0.30 on the registered slope. None changes the registered
verdict.

Two exploratory checks, pre-specified the same day as the interaction test, ask
how precise the SD 3.5 interval is. Letting each prompt have its own slope
leaves the estimates unchanged, -0.340 and -0.182, and widens the intervals, SD
3.5's to {[}-0.337, -0.027{]} and SDXL's to {[}-0.448, -0.231{]}, with
random-slope SDs of 0.112 (SDXL) and 0.177 (SD 3.5). Resampling prompts rather
than images widens the SD 3.5 slope interval to {[}-0.322, -0.050{]} and its
Spearman interval to {[}-0.407, -0.001{]}, while resampling seeds changes little
(slope {[}-0.238, -0.125{]}); with only six prompts these bootstrap intervals
are themselves imprecise.

\subsection{The distortion field and VLM-rated
veridicality}\label{the-distortion-field-and-vlm-rated-veridicality}

Exp 03b scores the identical images on axes adapted from Suzuki
\citep{suzuki2024}, pre-registered before any of those scores existed. The
registered claim required both checkpoints to pass three gates: a veridicality
slope of at least +0.20 with interval above zero, a partial correlation of
veridicality with guidance after adjustment for the Klüver composite of at least
+0.20, and kappa of at least 0.40. SDXL passes all three. SD 3.5 fails the first
two in the opposite direction: its veridicality slope is -0.261 {[}-0.347,
-0.174{]} and its partial is -0.161 {[}-0.262, -0.053{]}. Kappa on SD 3.5 is
0.485 and passes. The joint Exp 03b claim is not confirmed (Table 6). The SD 3.5
veridicality mixed model, like the primary composite model, did not converge
under L-BFGS; Powell did, at the variance boundary, with an estimate unchanged
to three decimals.

\begin{longtable}[]{@{}lllll@{}}
\caption{Pre-registered Exp 03b replication and dissociation gates on the
Suzuki-adapted axes. P1 is the guidance-standardised veridicality LMM slope
(registered pass \textgreater= +0.20, BH-adjusted p \textless= 0.05, CI above
zero); P2 is partial Spearman of veridicality and guidance adjusted for the
Klüver composite (pass \textgreater= +0.20, CI above zero); P3 is mean
field-level quadratic-weighted kappa (pass \textgreater= 0.40). Result is the
registered verdict.}\tabularnewline
\toprule\noalign{}
& P1 slope (95\% CI) & P2 partial rho (95\% CI) & P3 kappa & result \\
\midrule\noalign{}
\endfirsthead
\toprule\noalign{}
& P1 slope (95\% CI) & P2 partial rho (95\% CI) & P3 kappa & result \\
\midrule\noalign{}
\endhead
\bottomrule\noalign{}
\endlastfoot
SDXL & 0.463 {[}0.382, 0.543{]} & 0.417 {[}0.322, 0.506{]} & 0.529 & all pass \\*
SD 3.5 & -0.261 {[}-0.347, -0.174{]} & -0.161 {[}-0.262, -0.053{]} & 0.485 & P3
only \\*
\end{longtable}

The exploratory contrast, labelled as such, then asks how much of the distortion
field's movement survives adjustment for VLM-rated veridicality.

On SDXL, little of it. The correlation between the distortion field and guidance
falls from -0.486 raw to -0.175 {[}-0.27, -0.08{]} after adjustment for
VLM-rated veridicality, and per prompt from oranges, -0.70 to -0.08, through
living room, -0.53 to -0.385, to portrait, -0.23 to +0.24, where it changes sign
(Figure 4). On this exploratory measure the absolute correlation attenuated by
64\% after adjustment for veridicality, itself rated by the same judges. That is
construct overlap, not a causal split of sketchiness from warped objects. The
residual excludes zero, so the field is not empty. The judge screen gave one
concrete case of what the field bundles (Figure 5): on a painterly still life at
low guidance I scored condensation and distortion at 3 and Qwen3-VL-32B scored
both near 0. Neither was careless. The image was a loose painting of the
intended objects; I read the looseness as objects melting and the judge read it
as a style. That disagreement is the construct problem in one image, and it is
why the adjustment for veridicality was pre-registered for the second pass. Exp
01's primary endpoint was the form-constant metric, not this one; distortion was
its strongest exploratory field, and on SDXL most of that field's association
attenuated under the veridicality adjustment.

On SD 3.5 the same adjustment suppresses rather than attenuates: the raw -0.114
becomes -0.179 {[}-0.28, -0.08{]}, so the realism axis was masking the signal
rather than producing it, and a share-explained figure is undefined in that
direction.

Both patterns are possible on one field because distortion is bidirectional: one
ordinal absorbs under-conditioned melt at the bottom of the dial and
over-conditioned waxiness at the top. That construct problem is why the
composite without distortion is a sensitivity analysis in Table 5, not a
repaired primary endpoint.

The realism axis overshoots, exploratory: veridicality peaks at g = 7 on SD 3.5,
falling 1.07 rubric points by g = 15, and at g = 11 on SDXL, falling 0.16. The
top of the dial is not the most veridical setting.

\begin{figure}
\centering
\includegraphics[width=0.85\linewidth,height=\textheight,keepaspectratio,alt={The distortion field against VLM-rated veridicality, on SDXL. Per-prompt Spearman correlation between judge-mean distortion and guidance, raw (circles) against the partial correlation adjusting for veridicality (squares), sorted by the raw value. About 70 images per prompt; exploratory, uncorrected for multiple comparisons.}]{figures/fig3_style_confound.png}
\caption{The distortion field against VLM-rated veridicality, on SDXL.
Per-prompt Spearman correlation between judge-mean distortion and guidance, raw
(circles) against the partial correlation adjusting for veridicality (squares),
sorted by the raw value. About 70 images per prompt; exploratory, uncorrected
for multiple comparisons.}\label{fig:style}
\end{figure}

\begin{figure}
\centering
\includegraphics[width=0.55\linewidth,height=\textheight,keepaspectratio,alt={One image, two readings. SDXL, the still-life prompt, g = 1, seed 50, from the blind human subset. Scores on reduplication, fragmentation, condensation and distortion (0 to 3): the author 1, 2, 3, 3; Qwen3-VL-32B 0, 0, 0, 1; Claude Sonnet 5 2, 1, 1, 2. The intended watermelon, glass and keys are present but painted loosely, and the keys have dissolved into brushstrokes. Read as objects melting, it is high distortion and condensation; read as a painting style, it is close to none.}]{figures/fig5_painterly_stilllife.png}
\caption{One image, two readings. SDXL, the still-life prompt, g = 1, seed 50,
from the blind human subset. Scores on reduplication, fragmentation,
condensation and distortion (0 to 3): the author 1, 2, 3, 3; Qwen3-VL-32B 0, 0,
0, 1; Claude Sonnet 5 2, 1, 1, 2. The intended watermelon, glass and keys are
present but painted loosely, and the keys have dissolved into brushstrokes. Read
as objects melting, it is high distortion and condensation; read as a painting
style, it is close to none.}\label{fig:painterly}
\end{figure}

\subsection{Judge screening and silent
failure}\label{judge-screening-and-silent-failure}

Every candidate judge scored the same stratified subset of 28 images, on the
rubric it would use, under the blinding and decoding of the run. On the Klüver
rubric that screen was run on 2026-08-08, after the full pass of 2026-07-22 and
not before it; on the axes rubric it ran before the pass, as the rule the
programme adopted a day later requires. The screen asks one question: does the
judge use the scale. Two models did not use it at all. Qwen2.5-VL-7B and
Qwen3-VL-8B returned zero on every field of every one of the 28 images, so their
fraction of images scored above zero is flat at the floor for reduplication,
fragmentation, condensation and distortion alike (Figure 6).
Llama-3.2-11B-Vision scored above zero on reduplication for 18\% of the subset
and was flat at zero on the other three fields. Gemma-3-27B used two of the
four, reduplication on 36\% of the subset and distortion on 29\%, and was flat
on fragmentation and condensation. Qwen3-VL-32B and Claude Sonnet 5 were the two
that used the scale, and they used it differently. On the probe the 32B moved on
reduplication and distortion, at 25\% of the subset each, on condensation for
11\%, and on fragmentation for one image in 28, 4\%. It passes the screen, which
asks only for non-zero variance on every graded field, tiling excluded; its own
tiling column is flat and the screen does not read it. Its profile is not
Claude's.

Nothing about the two silent-zero runs looked wrong. For Qwen2.5-VL-7B and
Qwen3-VL-8B there were no refusals, no decoding errors and no schema violations:
every reply was well-formed JSON, and each carried a fluent, confident free-text
note about the image. Llama-3.2-11B-Vision failed in the way a pipeline does
catch, and is not part of this claim: roughly 90 of its replies were truncated
JSON at a max\_tokens of 400, and the completed ones were flat at zero on three
of the four fields. On a low-guidance SDXL living room the 7B judge wrote ``The
image depicts a coherent living room with one sofa, one lamp, and one window,
without any reduplication, fragmentation, condensation, or distortion'' and
scored the row zero throughout. A pipeline that monitors parse rates, refusal
rates and schema validity would have passed the two Qwen judges and reported
their constant columns as a fact about the images.

The repair was a post-data instrument amendment, dated 2026-08-08 and taken
under the pre-registration clause that lists judge models as swappable without
reclassifying the experiment, on the trigger condition the pre-registration had
named, a poor human-subset kappa for the 7B, which turned out to be zero. It was
not specified in advance as a swap to this model. The pre-registration named
Qwen2.5-VL-32B as the fallback; the model that cleared the screen and took the
seat is Qwen3-VL-32B, one generation newer, at the same quantisation, and the
amendment records that substitution as a deviation. Nothing else moved: the same
rubric, the same prompt, the same decoding settings and the same images.
Composite (mean per-field) quadratic-weighted kappa between the two remaining
judges rose from 0.290 to 0.562 on SDXL and from 0.158 to 0.440 on SD 3.5 (Table
7). The panel as first run, which added Llama-3.2-11B-Vision, was lower still at
0.126 and 0.135, because a rater whose column never varies pulls a weighted
kappa toward zero by arithmetic rather than by disagreement. The original run is
retained in the repository beside the amended one; it is not replaced.

\begin{longtable}[]{@{}
  >{\raggedright\arraybackslash}p{(\linewidth - 4\tabcolsep) * \real{0.3333}}
  >{\raggedright\arraybackslash}p{(\linewidth - 4\tabcolsep) * \real{0.3333}}
  >{\raggedright\arraybackslash}p{(\linewidth - 4\tabcolsep) * \real{0.3333}}@{}}
\caption{Composite (mean per-field) quadratic-weighted kappa across judge panels
on the same generated corpus and rubric. Row-specific n is the complete-case
image count and differs because Llama had missing replies. The first two rows
are the panels as run on 2026-07-22; the third is the amended panel, after the
Qwen2.5-VL-7B seat was refilled with Qwen3-VL-32B on 2026-08-08, with no rubric
change. Human rows use a 28-image blind subset drawn from both checkpoints, too
few to split, so each row has one pooled value, shown under SDXL; bootstrap 95\%
CI. The second-rater rows are exploratory, with no gate.}\tabularnewline
\toprule\noalign{}
\begin{minipage}[b]{\linewidth}\raggedright
panel
\end{minipage} & \begin{minipage}[b]{\linewidth}\raggedright
SDXL
\end{minipage} & \begin{minipage}[b]{\linewidth}\raggedright
SD 3.5
\end{minipage} \\
\midrule\noalign{}
\endfirsthead
\toprule\noalign{}
\begin{minipage}[b]{\linewidth}\raggedright
panel
\end{minipage} & \begin{minipage}[b]{\linewidth}\raggedright
SDXL
\end{minipage} & \begin{minipage}[b]{\linewidth}\raggedright
SD 3.5
\end{minipage} \\
\midrule\noalign{}
\endhead
\bottomrule\noalign{}
\endlastfoot
Claude + Qwen2.5-VL-7B (original judge B) & 0.290 (n = 428) & 0.158 (n = 427) \\*
Claude + Qwen2.5-VL-7B + Llama-3.2-11B (original three-judge panel) & 0.126 (n =
380) & 0.135 (n = 387) \\*
Claude + Qwen3-VL-32B (amended panel) & 0.562 (n = 428) & 0.440 (n = 427) \\*
human (author, n = 28) vs Claude Sonnet 5 & 0.337 {[}0.106, 0.495{]} &
(pooled) \\*
human (author, n = 28) vs Qwen3-VL-32B & 0.124 {[}-0.059, 0.313{]} & (pooled) \\*
second rater (naive, n = 28) vs Claude Sonnet 5 & 0.426 {[}0.208, 0.575{]} &
(pooled) \\*
second rater (naive, n = 28) vs Qwen3-VL-32B & 0.169 {[}-0.029, 0.391{]} &
(pooled) \\*
author vs second rater (human vs human, n = 28) & 0.567 {[}0.330, 0.721{]} &
(pooled) \\*
\end{longtable}

Agreement between models is not validity. The first rater, the author, scored a
blind subset of 28 images with guidance and architecture hidden and the order
shuffled. Against that rater, Claude Sonnet 5 reaches a composite (mean
per-field) kappa of 0.337 {[}0.106, 0.495{]} and Qwen3-VL-32B reaches 0.124
{[}-0.059, 0.313{]}, an interval that spans zero. On those same 28 images the
two judges agree with each other at a mean per-field kappa of 0.271 {[}0.056,
0.448{]}, below human against Claude at 0.337 {[}0.106, 0.495{]} and above human
against Qwen at 0.124. The full-corpus panel kappa of 0.562 and 0.440 uses n =
428 and n = 427 images, including the ten empty-prompt baselines per checkpoint;
these subset values are computed on different image sets and are not compared
directly. Per field on the 28, the judges agree least on distortion, at 0.069,
the field where human against Claude agreement is also weakest. The weakest
field for the better of them is distortion, at 0.212 human against Claude. Exp
01's primary endpoint was the form-constant metric; distortion was its strongest
exploratory field.

A second rater, who does not know the hypothesis, later scored the same 28
images in the same order, from a plain-language version of the rubric
(exploratory, no gate; the wording is in the repository). The two humans agree
at a composite kappa of 0.567 {[}0.330, 0.721{]}, about as well as the two
judges agree with each other on the full SDXL corpus. Against the second rater,
Claude Sonnet 5 reaches 0.426 {[}0.208, 0.575{]}, higher than against me, so the
author's knowledge of the hypothesis did not inflate the judge's human
agreement; Qwen3-VL-32B reaches 0.169 {[}-0.029, 0.391{]}, an interval that
again spans zero. With 28 images every one of these intervals is wide.

I read this as three tiers of qualification for a judge on a breakdown rubric:
it uses the scale on a stratified probe; it agrees with a judge from another
model family, which on the full corpus this panel does at 0.562 and 0.440,
although on the 28-image human subset the two judges agree with each other less
than Claude agrees with the rater; and it agrees with a human rater. The panel
reported here clears the first two and clears the third only partly. A judge can
pass the rubric screen and still lack human validity.

The size reading of this result is narrow. Five open-weight judges were
screened, and the two that returned no variance at all belong to a single Qwen
lineage across two generations, so parameter count is not separable here from
model family or release date. ``Below about 30B'' describes this panel and not a
law, and Gemma-3-27B's partial result, graded on two fields and flat on two,
already blurs any threshold. The direction is consistent with Chen et al., who
report that MLLM judges track human preference on pair comparison while
diverging from it on scoring evaluation, that the open-weight models they
evaluate (LLaVA, CogVLM, Qwen-VL-Max) align worse with human judgement than
GPT-4V, and that judge outputs carry biases and hallucinated justifications
\citep{chen2024mllmjudge}. What I observe is more extreme than divergence: no
variance to diverge with.

The scored images (835 conditioned and 20 empty-prompt images with two complete
judge replies, out of 860 generated), the raw replies behind them and the
28-image human subset are public, and can be reused as a screening set for other
judges.

\begin{figure}
\centering
\includegraphics[width=0.85\linewidth,height=\textheight,keepaspectratio,alt={Judge screen. Per cell, the percentage of the 28 blind-subset images each rater scored above zero (upper number) and the mean score on the 0 to 3 scale (lower number), per field, on the shared rubric; parse failures are excluded from the denominator. Qwen2.5-VL-7B and Qwen3-VL-8B sit at zero on every field, Llama-3.2-11B-Vision moves on reduplication only, Gemma-3-27B on reduplication and distortion only, and Qwen3-VL-32B and Claude Sonnet 5 score above zero on all four, at rates that differ between them. The screen reads the four graded fields only; the binary tiling flag is excluded, so Qwen3-VL-32B passes although its tiling column is flat. The human row is the single rater, the author, on the same images.}]{figures/fig4_judge_screen.png}
\caption{Judge screen. Per cell, the percentage of the 28 blind-subset images
each rater scored above zero (upper number) and the mean score on the 0 to 3
scale (lower number), per field, on the shared rubric; parse failures are
excluded from the denominator. Qwen2.5-VL-7B and Qwen3-VL-8B sit at zero on
every field, Llama-3.2-11B-Vision moves on reduplication only, Gemma-3-27B on
reduplication and distortion only, and Qwen3-VL-32B and Claude Sonnet 5 score
above zero on all four, at rates that differ between them. The screen reads the
four graded fields only; the binary tiling flag is excluded, so Qwen3-VL-32B
passes although its tiling column is flat. The human row is the single rater,
the author, on the same images.}\label{fig:judges}
\end{figure}

\subsection{Form constants: the registered validation failed on
specificity}\label{form-constants-the-registered-validation-failed-on-specificity}

Exp 01 asked the form-constant question directly. One judge, Claude Sonnet 4.6,
flagged Klüver's lattice, cobweb, tunnel and spiral \citep{kluver1966} and
scored a geometric intensity, the pre-registered metric M (0 to 1). M sat near
the floor at every guidance value on both checkpoints, with per-guidance means
between 0.033 and 0.167. The registered Spearman of those nine means against
guidance is -0.044 on SDXL and +0.234 on SD 3.5 (+0.01 and +0.10 across the 90
images). Both readings sit inside the pre-registered null region of
\textbar rho\textbar{} \textless{} 0.3; confirmation required rho at or below
-0.6. What that null is worth depends entirely on whether the judge could see a
form constant at all, and after Section 3.4 a flat reading from an unchecked
judge is worth little.

Exp 02 Stage A is that check, pre-registered before any image was scored. The
test images come from a simplified model proposed to generate form-constant
patterns, run as code: a minimal Ermentrout-Cowan field on a periodic cortical
sheet, pushed through its Turing bifurcation and rendered to the visual field
through the inverse log-polar map \citep{ermentrout1979}. Three of its four
gates pass: M is 1.00 at every gain above threshold and 0.00 on the blank
renders below it, all four classes appear, and M rises with gain (Spearman 0.87
hexagon, 0.91 stripe). The fourth, which required the cortical pattern to
predict the judged class, fails on the tunnel and lattice rows and passes on the
spiral.

\begin{longtable}[]{@{}lll@{}}
\caption{Exp 02 positive control for the Exp 01 form-constant judge (archived
rubric, single blind call per image). A set counts as flagged when any of the
four binary class flags fires. Pre-specified pass required the ordinary-image
rate at or below 20\%.}\tabularnewline
\toprule\noalign{}
image set & flagged & rate (95\% CI) \\
\midrule\noalign{}
\endfirsthead
\toprule\noalign{}
image set & flagged & rate (95\% CI) \\
\midrule\noalign{}
\endhead
\bottomrule\noalign{}
\endlastfoot
rendered form constants (mu \textgreater= 1.05 mu\_c) & 80/80 & 100\% {[}95\%,
100\%{]} \\*
blank renders (mu \textless= 0.9 mu\_c) & 0/40 & 0\% {[}0\%, 9\%{]} \\*
ordinary Exp 03 images (g in \{7, 11\}) & 18/40 & 45\% {[}31\%, 60\%{]} \\*
\end{longtable}

Table 8 gives the judge test. Sensitivity is perfect: all 80 rendered form
constants are flagged and no blank is. Specificity is not: the flags also fire
on 18 of 40 ordinary generated images from the Exp 03 grid at g in \{7, 11\},
45\% against a registered ceiling of 20\%. The false positives are not invented
geometry; every one names a literal grid in its note, a brick wall, window
mullions, tiles, and the portrait prompt, which contains no grid, was flagged 0
of 8. The pre-registration fixed what that means: the instrument is not
validated, and the Exp 01 null is reclassified as uninterpretable. It is not a
finding about diffusion output.

Three limits attach. The judge is not class-specific: under the log-polar map
the hexagonal lattice renders with a global spiral arrangement, so lattice
renders draw the spiral and tunnel flags too. The registered criterion used any
flag rather than M, the wrong summary statistic; exploratory, stated after the
table was seen, on M itself every rendered form constant reaches the top score
and no ordinary image does. And the simulated curve is threshold-shaped, flat
below the bifurcation and saturated just above it; no test of correspondence
with the guidance curves was pre-registered, so none is claimed.

\section{Discussion}\label{discussion}

\subsection{What this paper does not
claim}\label{what-this-paper-does-not-claim}

This paper does not claim that lowering classifier-free guidance models
psychedelic prior relaxation, that form constants appear in diffusion output,
that any operation shown here transfers to a brain, or that the effect is
specific to guidance rather than to image degradation in general; and it does
not call SD 3.5 a null, because its linear effect is reliably non-zero and its
interval includes the registered gate of -0.20 rather than excluding an effect
of that size.

The mapping between the guidance dial and the psychedelic-neuroscience
literature is a modelling choice, and both readings set out in Section 1 remain
open: lowering guidance weakens a top-down prior if the prompt is read as the
prior, and increases prior dominance if the model's unconditional distribution
is. Nothing in my design discriminates between them. The neuroscience supplies
the rubric and the perturbation framing; it is not tested here.

I perturbed one knob. No alternative degradation, such as fewer sampling steps,
injected noise, or prompt ablation, was scored on the same rubric, so I cannot
claim that the profile reported above is specific to guidance rather than a
generic signature of image degradation. The quality adjustment uses two
no-reference quality proxies; it is not a second perturbation.

The registered validation of the form-constant instrument failed on specificity,
so under the rule fixed in advance the Exp 01 null is reclassified as
uninterpretable: the archived judge scores every rendered form constant at the
top of its scale and every blank at the bottom, but its binary class flags also
fire on ordinary generated scenes that contain a literal grid, above the ceiling
registered for that rate. The sensitivity reading is therefore exploratory, and
the Exp 01 null is neither a finding about geometric structure nor an
established absence of one.

No biological comparison arm exists in this program. I collected no human
hallucinatory or neurobiological comparison data and computed no reference
curve, so nothing here compares diffusion breakdown against a biological
measurement.

On SD 3.5 the pre-registered linear slope excludes zero, and its interval
includes the registered gate of -0.20 on the pre-registered scale, so the
criterion is not met and an effect of the registered size is not excluded
either; the effect is carried by the g = 1 condition. I report a reliable effect
that misses the registered bar, not an absence of one.

The contrast between a graded profile on one checkpoint and a change
concentrated at the bottom of the dial on the other is post-hoc, and is labelled
as such wherever it appears. The exploratory sweep below the lowest registered
guidance value is not part of this paper.

\subsection{What the assay is for}\label{what-the-assay-is-for}

The assay is a package another group can pick up and run. It is a doseable
perturbation on a knob that most text-to-image models already expose, a rubric
whose fields name phenomena rather than rate quality, a screen that disqualifies
a judge before a run is spent on it, and a public set of scored images, with the
raw judge replies and the blind human subset beside them, usable as a screening
set for any other judge on either rubric. The same rubric and images can serve
the other way round, as a benchmark for generative models: a new checkpoint run
on the six prompts at the seven guidance values can be scored by the amended
panel and placed against these two curves. That is a use of the assay, not a
result of it. None of that depends on the neuroscience that supplied the rubric.

\subsection{Future work}\label{future-work}

These are the follow-ups I would run next, in order, and that anyone can run on
the public images and code.

\begin{enumerate}
\def\labelenumi{\arabic{enumi}.}
\tightlist
\item
  \textbf{Specificity of the perturbation.} No alternative perturbation was
  scored on the same rubric: not fewer sampling steps, not injected noise, not
  timestep-restricted guidance, not prompt ablation. Until one is, I cannot say
  the profile reported here belongs to guidance rather than to image degradation
  in general. It is also the explicit reason I do not claim the analogy, and the
  next pre-registration is that comparison.
\item
  \textbf{A human rater panel.} One naive second rater on 28 images (Section
  3.4) shows the rubric is visible to people; several raters on a larger subset,
  with the judges' exact rubric wording, are what would turn the endpoint into a
  measure of human-perceived breakdown.
\item
  \textbf{A second checkpoint per architecture.} Without it the difference
  between the two checkpoints cannot be attributed to the architecture.
\item
  \textbf{Prompt adherence as a covariate.} Guidance is the prompt-following
  dial, so a prompt-image alignment score, measured and adjusted for next to the
  quality proxies, would test whether breakdown is prompt-following under
  another name.
\item
  \textbf{A shape-agnostic endpoint.} The registered straight-line slope
  understated an effect concentrated in one step (Section 3.2); the next
  registration should not assume a shape.
\end{enumerate}

Two other readings of the same knob, a therapeutic reading of a guidance
schedule and an annealing reading of the sampler, appear in the public notebook.
They are framing, not results here.

\subsection{What I learned about doing
this}\label{what-i-learned-about-doing-this}

Three rules entered the methodology halfway through the programme, and each was
paid for by a failure inside it.

Do not let the endpoint assume the shape of the response. A standardised linear
slope on a grid that is close to geometric could not tell a gradient from a
change concentrated at the bottom of the dial, and reported the two checkpoints
as one effect at two sizes. A monotone trend test or an explicit segmented fit
was the right registration and is what I would register now.

Never quote a single pooled correlation. The per-prompt range is the finding. On
SDXL the per-prompt correlations of composite against guidance run from -0.74 to
-0.28, and the attenuation of the distortion association after the veridicality
adjustment runs from almost all of it on one prompt to under a third on another.
A single number averages over that spread and is not the effect.

Screen the judge on the rubric it will actually score, before the full run. Two
of my judges failed silently, and I only found out from the human subset after
the run was scored. In Exp 03 the screen came after the run, and the panel
change is therefore a post-data amendment. The rule is what this experiment
produced, not what it followed.

\section{Limitations}\label{limitations}

The first human rater is me, and I know the hypothesis. Guidance and
architecture were hidden from the rater and the order was shuffled, but a g = 1
SDXL image is visually identifiable as low-guidance, so the rater can infer the
condition from the stimulus. ``Blind to guidance'' is weaker than it sounds.
Human against Qwen3-VL-32B composite kappa is 0.124 with an interval spanning
zero, so the second judge has no demonstrated human validity at all. The panel's
0.562 and 0.440 are model-model agreement on the full corpus; on the same 28
images the two judges agree at 0.271, below my own 0.337 with Claude. A second
rater who does not know the hypothesis agrees with me at 0.567 and with Claude
at 0.426, which shows the rubric describes something two people can see; but two
raters on 28 images, one working from a plain-language version of the rubric,
are not a validated measure of human-perceived breakdown.

The panel carries fewer judge lineages than the design budgeted: two model
families in the final panel, five open judges screened, and one Qwen lineage
across two generations, so parameter count is not separable here from family or
release.

The distortion field mixes two opposite failure modes, under-conditioned melt
and over-conditioned waxiness, on one ordinal. And 64\% of its absolute
correlation with guidance on SDXL attenuated after adjustment for veridicality
rated by the same judges. That 64\% is construct overlap, not a causal
decomposition. The three-field composite is a sensitivity analysis, not a
repaired endpoint.

SD 3.5's effect rests on g = 1, the single most under-conditioned and therefore
most confound-prone level of the dial, and its interval includes the registered
gate of -0.20. A prompt-cluster bootstrap of the quality-adjusted partial
correlation, pre-specified on 2026-09-25 as sensitivity, gives {[}-0.570,
-0.251{]} on SDXL and {[}-0.410, 0.024{]} on SD 3.5 around points of -0.433 and
-0.245. It does not change the registered verdict, and it shows that on SD 3.5
the quality-adjusted association is not robust to which prompts were run. The
two quality components, CLIP-IQA and the aesthetic predictor, agree poorly with
each other, Spearman -0.155 on SDXL and +0.170 on SD 3.5; adjusting for either
component separately still leaves a negative partial correlation.

The registered SD 3.5 kappa of 0.440 includes the ten empty-prompt baselines. On
conditioned images only it is 0.394, which misses the 0.4 gate. The SD 3.5 mixed
model and the Exp 03b veridicality mixed model both failed to converge under
L-BFGS and were refit with Powell at the variance boundary; the slopes did not
change at the reported precision.

The registered Exp 03b claim fails: SDXL passes P1, P2 and P3, and SD 3.5 fails
P1 and P2 in the opposite direction. The exploratory veridicality-adjusted
distortion contrast does not replace that registered result.

The pre-registered linear endpoint on a near-geometric grid could not
distinguish a gradient from a change concentrated at g = 1. That contrast is
post-hoc.

The forest control's composite tracked guidance on SDXL, so the object-bound
reading of the breakdown is not established.

The grids differ between Exp 01 and Exp 03, so the two are not pooled. Exp 02 is
single-judge and single-call with no human check, and its registered control
failed on specificity. Both are latent text-to-image checkpoints from one vendor
family, one a UNet denoiser and one a rectified-flow transformer, one checkpoint
each, so architecture here is confounded with training data, training objective,
text encoder and sampler.

The judge swap of 2026-08-08 was post-data. The pre-registered fallback was
Qwen2.5-VL-32B, and Qwen3-VL-32B was used in its place as a recorded deviation.

The calibration of Sonnet 5 against Exp 01's Sonnet 4.6 was pre-registered and
not run.

Image-text similarity was not measured. Prompt adherence rises with guidance in
general, so a reader may ask whether the breakdown score is prompt-following
under another name. The quality adjustment used two no-reference scores and set
CLIP similarity aside because it measures obedience rather than image quality;
that leaves prompt adherence as an unmeasured competing explanation of the
composite, and the next pre-registration should score it.

\section{Reproducibility}\label{reproducibility-1}

Code, pre-registrations, analyses and all judge outputs are at
\url{https://github.com/youssefhassan/operating-system-hypothesis-public}.
Generated images and judgements are published as datasets under
\texttt{youssefhassan13/} on Hugging Face, named exp01-guidance-sweep,
exp02-form-constant-generator and exp03-l23-hardening.

\begin{verbatim}
# Exp 01
cd experiments/01_image_test
python analyze.py --model sdxl --plots && python analyze.py --model sd35 --plots

# Exp 02
cd experiments/02_form_constant_generator && python analyze.py

# Exp 03, primary (amended two-judge panel)
cd experiments/03_l23_hardening
python analyze.py --both --plot --judges claude,qwen

# Exp 03, post-hoc and both sensitivity rounds
python posthoc.py --judges claude,qwen
python posthoc_sensitivity.py --judges claude,qwen
python posthoc_sensitivity2.py --judges claude,qwen
python posthoc_quality_paths.py --judges claude,qwen

# Exp 03b, Suzuki axes
python analyze_axes.py --judges claude,qwen
\end{verbatim}

Each command regenerates its report JSON from the public repository and the
published datasets. Byte-for-byte regeneration was verified on 2026-09-16 and
2026-09-17; the analysis code was corrected on 2026-09-29 (dated notes in each
experiment's analysis.md), and the corrected reports have not yet been
re-verified from a clean public checkout. The pre-registration dates quoted
throughout are the git commit dates of the pre-registration files on that public
repository.

\section{Use of AI tools}\label{use-of-ai-tools}

Two vision-language models, Claude Sonnet 5 and Qwen3-VL-32B, are instruments of
this study: they score the images, and their screening and failures are reported
in Section 2.5 and Section 3.4. Separately, I used Claude (Anthropic) as a
coding and writing assistant. Under my direction it wrote and ran analysis and
figure code, ran the automated checks that trace every number in the text to an
analysis output, verified every reference against its source, and drafted and
edited passages of the text. I designed the study, made every methodological and
interpretive decision, reviewed every claim against the data, and take full
responsibility for the content.

\section{Appendix: Figure 2 as numbers}\label{appendix-figure-2-as-numbers}

The fourteen means behind each line of Figure 2, so the figure can be re-plotted
without the repository. Both columns per checkpoint are on that checkpoint's own
standardised ruler and are not compared across checkpoints.

\begin{longtable}[]{@{}
  >{\raggedright\arraybackslash}p{(\linewidth - 12\tabcolsep) * \real{0.1429}}
  >{\raggedright\arraybackslash}p{(\linewidth - 12\tabcolsep) * \real{0.1429}}
  >{\raggedright\arraybackslash}p{(\linewidth - 12\tabcolsep) * \real{0.1429}}
  >{\raggedright\arraybackslash}p{(\linewidth - 12\tabcolsep) * \real{0.1429}}
  >{\raggedright\arraybackslash}p{(\linewidth - 12\tabcolsep) * \real{0.1429}}
  >{\raggedright\arraybackslash}p{(\linewidth - 12\tabcolsep) * \real{0.1429}}
  >{\raggedright\arraybackslash}p{(\linewidth - 12\tabcolsep) * \real{0.1429}}@{}}
\caption{The two lines of Figure 2 as numbers: complete-case n, mean breakdown
composite and mean quality score Q at each guidance value, in standardised units
within each checkpoint. Conditioned cells contain 58 to 60 images after parser
failures; the empty-prompt row has no quality score in the pre-registered
analysis.}\tabularnewline
\toprule\noalign{}
\begin{minipage}[b]{\linewidth}\raggedright
g
\end{minipage} & \begin{minipage}[b]{\linewidth}\raggedright
SDXL n
\end{minipage} & \begin{minipage}[b]{\linewidth}\raggedright
SDXL composite
\end{minipage} & \begin{minipage}[b]{\linewidth}\raggedright
SDXL Q
\end{minipage} & \begin{minipage}[b]{\linewidth}\raggedright
SD 3.5 n
\end{minipage} & \begin{minipage}[b]{\linewidth}\raggedright
SD 3.5 composite
\end{minipage} & \begin{minipage}[b]{\linewidth}\raggedright
SD 3.5 Q
\end{minipage} \\
\midrule\noalign{}
\endfirsthead
\toprule\noalign{}
\begin{minipage}[b]{\linewidth}\raggedright
g
\end{minipage} & \begin{minipage}[b]{\linewidth}\raggedright
SDXL n
\end{minipage} & \begin{minipage}[b]{\linewidth}\raggedright
SDXL composite
\end{minipage} & \begin{minipage}[b]{\linewidth}\raggedright
SDXL Q
\end{minipage} & \begin{minipage}[b]{\linewidth}\raggedright
SD 3.5 n
\end{minipage} & \begin{minipage}[b]{\linewidth}\raggedright
SD 3.5 composite
\end{minipage} & \begin{minipage}[b]{\linewidth}\raggedright
SD 3.5 Q
\end{minipage} \\
\midrule\noalign{}
\endhead
\bottomrule\noalign{}
\endlastfoot
1 & 60 & 1.108 & -0.691 & 60 & 1.094 & -0.591 \\*
2 & 60 & 0.224 & -0.127 & 58 & -0.151 & -0.153 \\*
3 & 60 & -0.051 & -0.112 & 60 & -0.166 & -0.008 \\*
5 & 59 & -0.242 & 0.049 & 60 & -0.260 & 0.118 \\*
7 & 59 & -0.360 & 0.156 & 59 & -0.280 & 0.312 \\*
11 & 60 & -0.340 & 0.291 & 60 & -0.165 & 0.364 \\*
15 & 60 & -0.350 & 0.438 & 60 & -0.082 & -0.043 \\*
empty prompt & 10 & 1.510 & n/a & 10 & 2.172 & n/a \\*
\end{longtable}

\section{Appendix: the judge prompts}\label{appendix-the-judge-prompts}

Each judge receives the image and one of the two prompts below, verbatim. The
first sentence of each names the intended scene; it is shown here for the
still-life prompt (``a watermelon, a glass half-filled with water, and a set of
keys on a wooden table''). The guidance value never appears. Source:
\texttt{rubric.py} and \texttt{rubric\_axes.py} in the Exp 03 directory of the
repository.

\textbf{Klüver rubric (Exp 03).}

\begin{verbatim}
Score these phenomena from Heinrich Klüver's Level-2 and Level-3
hallucinatory
constants. This image was meant to depict: "a watermelon, a glass
half-filled with
water, and a set of keys on a wooden table". Intended content: one
watermelon, one
glass, one set of keys on a tabletop. Judge only what is VISIBLY
present; do not guess
what it was supposed to look like. Reduplication means repetition
BEYOND that intended
content.

Score:
- reduplication: objects or motifs repeated far beyond the
  intended content
  (e.g. many copies where one was intended, rows/fields of an
  object). 0 none,
  3 dominant.
- fragmentation: objects broken into disconnected or shattered
  pieces. 0 none,
  3 dominant.
- condensation: distinct objects fused/merged into hybrid forms
  (chimeras, one
  object melting into another). 0 none, 3 dominant.
- distortion: melted, warped, or anatomically/physically
  impossible shapes
  (dysmorphopsia). 0 none, 3 dominant.
- tiling: 1 if the whole image has dissolved into a repeating
  texture / pattern
  field rather than a single readable scene, else 0.

Return STRICT JSON only, no prose, with exactly these keys:
{"reduplication": 0..3, "fragmentation": 0..3, "condensation":
0..3,
 "distortion": 0..3, "tiling": 0 or 1, "notes": "one short
 phrase"}
\end{verbatim}

\textbf{Suzuki-adapted axes (Exp 03b).} The field definitions, verbatim from Exp
01:

\begin{verbatim}
- veridicality: 0 abstract/unreal, 3 photoreal coherent scene.
- spontaneity: 0 fully prompt-driven, 3 content the prompt didn't ask for.
- complexity: 0 simple, 3 highly elaborate.
- coherent_scene: 1 if the objects hang together as one believable scene, else 0.
\end{verbatim}

Line breaks inside the prompts are added here for the page width; the text is
unchanged. The ``Level'' numbering in the prompt's first line is this project's
shorthand, not Klüver's.
\bibliographystyle{unsrtnat}
\bibliography{refs}
\end{document}
